% =========================================================================
%  Ejercicio 2b - Ensayo INDIVIDUAL sobre el articulo
%  Responsable: Juárez Larrondo Omar Alejandro
% =========================================================================

\newpage
\subsection*{Ejercicio 2b. Ensayo individual}

\begin{flushright}
\textbf{Autor: Juárez Larrondo Omar Alejandro}
\end{flushright}

\subsubsection*{Introducción}

La lectura de \citet{datavalue} aborda una realidad que suele permanecer oculta
tras la comodidad de los servicios digitales: los datos personales poseen un
valor económico y constituyen la materia prima de numerosos modelos de negocio.
El tema resulta relevante porque la vida cotidiana depende cada vez más de
plataformas, teléfonos, sistemas de recomendación y objetos conectados que
registran decisiones y comportamientos. Frente a este escenario, se sostiene
que el aprovechamiento de los datos puede producir beneficios sociales y
tecnológicos, pero solo es legítimo cuando existe transparencia, una finalidad
definida y capacidad efectiva de control por parte de sus titulares. La simple
aceptación de condiciones extensas no basta para equilibrar la relación entre
quien genera la información y quien obtiene beneficios de ella.

\subsubsection*{Desarrollo}

La autora explica que el Big Data y el Internet de las cosas intensifican la
recolección de información. Las redes sociales, los motores de búsqueda y los
sensores transforman actividades ordinarias en registros que pueden agregarse
y analizarse. El valor no reside necesariamente en un dato aislado, sino en la
capacidad de combinar millones de observaciones para reconocer patrones,
construir perfiles y anticipar conductas. De este modo, las empresas convierten
la información en publicidad dirigida, conocimiento del mercado y servicios
personalizados. El objetivo del artículo consiste en hacer visible este proceso
y persuadir sobre la necesidad de conocer tanto el valor monetario de los datos
como las consecuencias de ceder su control.

Esta postura resulta convincente porque el intercambio entre servicios y datos
no ocurre en condiciones plenamente simétricas. Una persona puede revelar su
ubicación para obtener una ruta, permitir el registro de reproducciones para
recibir recomendaciones o compartir su historial de compra para agilizar una
transacción. Cada finalidad parece limitada; sin embargo, la unión de esos
registros puede revelar horarios, nivel socioeconómico, relaciones y
preferencias que nunca fueron declaradas de manera directa. Por lo tanto, el
riesgo no se reduce al robo de una contraseña. También comprende inferencias
incorrectas, discriminación automatizada, manipulación comercial y usos
secundarios incompatibles con el propósito original.

La relación con Fundamentos de Bases de Datos es directa. Una base de datos no
es un recipiente neutral, pues su modelo determina qué aspectos de una persona
se representan, cómo se relacionan y durante cuánto tiempo permanecen
disponibles. Las decisiones sobre claves, restricciones de integridad,
permisos, bitácoras, respaldos y políticas de eliminación tienen consecuencias
para la privacidad. Asimismo, la calidad de los datos es indispensable: cuando
un perfil contiene errores o carece de contexto, una consulta correcta puede
producir una decisión injusta. En consecuencia, el diseño técnico debe incorporar
minimización, control de acceso, trazabilidad y separación de finalidades desde
el inicio, y no como correcciones posteriores.

La temática central de la lectura es la valorización económica de la información
personal. A su alrededor aparecen temas laterales inseparables: privacidad,
consentimiento, confianza, propiedad, portabilidad, interoperabilidad y
gobernanza. En la práctica profesional relacionada con sistemas de información,
estos asuntos obligan a preguntar si cada atributo solicitado es realmente
necesario, quién podrá consultarlo y qué procedimiento permitirá corregirlo o
eliminarlo. Por ejemplo, una aplicación que almacena ubicación para una función
momentánea no debería conservar indefinidamente el historial ni reutilizarlo
para publicidad sin una autorización comprensible y específica.

No obstante, reducir la solución a pagar a cada individuo por sus datos sería
insuficiente. Un precio no elimina el desequilibrio de conocimiento ni evita
que una necesidad económica conduzca a aceptar condiciones perjudiciales. La
propuesta más valiosa del texto es el empoderamiento mediante herramientas
centradas en la persona, como los almacenes personales de datos, junto con
derechos y obligaciones verificables. Tal enfoque permitiría decidir qué se
comparte y revocar permisos, a la vez que conservaría las oportunidades de
investigación e innovación. Su viabilidad dependería de interfaces claras,
estándares interoperables y auditorías que traduzcan los principios jurídicos en
controles técnicos comprobables.

\subsubsection*{Conclusión}

En conclusión, el uso de datos personales resulta aceptable cuando las personas
comprenden el intercambio y mantienen una capacidad real de decisión. La
lectura permite reconocer que el valor de la información no se limita a una
cifra: también incluye su poder para describir, predecir e influir sobre una
persona. Su aportación principal consiste en vincular la conciencia individual
con la responsabilidad de empresas, desarrolladores y autoridades. Quedan
abiertos retos como valorar datos derivados, explicar inferencias algorítmicas,
garantizar la eliminación en sistemas distribuidos y evitar que la protección de
la privacidad dependa del poder adquisitivo. Afrontarlos exigirá bases de datos
seguras y bien gobernadas, pero también una cultura profesional que considere la
información personal como una extensión de la dignidad y la autonomía humanas.
